% =========================================================
% Inequality: Algebra of Strict and Nonstrict Inequalities
% =========================================================

\subsection*{Inequality}
\label{sec:inequality}

% ---------------------------------------------------------
% thm:real-order-trichotomy
% ---------------------------------------------------------

\begin{theorem}[Trichotomy for Real Order]
\label{thm:real-order-trichotomy}
For all $a,b\in\mathbb{R}$, exactly one of the following relations holds:
\[
a<b,\qquad a=b,\qquad a>b.
\]
\hyperref[prf:real-order-trichotomy]{\textit{Go to proof.}}
\end{theorem}

\begin{remark*}[Standard quantified statement]
\[
\forall a,b\in\mathbb{R}\;
\bigl((a<b)\oplus(a=b)\oplus(a>b)\bigr).
\]
\end{remark*}

\begin{remark*}[Predicate reading]
\[
a\text{ and }b\text{ are comparable by exactly one real-order relation.}
\]
\end{remark*}

\begin{remark*}[Interpretation]
Any two real numbers can be compared, and the three possible comparison
outcomes are mutually exclusive.
\end{remark*}

\begin{dependencies}
\hyperref[thm:trichotomy-on-r]{Trichotomy on $\mathbb{R}$}.
\end{dependencies}

% ---------------------------------------------------------
% thm:ineq-add-both-sides
% ---------------------------------------------------------

\begin{theorem}[Addition Preserves Strict Inequality]
\label{thm:ineq-add-both-sides}
For all $a, b, c \in \mathbb{R}$,
\[
  a < b \;\Rightarrow\; a + c < b + c.
\]
\hyperref[prf:ineq-add-both-sides]{\textit{Go to proof.}}
\end{theorem}



\begin{remark*}[Standard quantified statement]
\[
  \forall a, b, c \in \mathbb{R},\quad
  a < b \;\Rightarrow\; a + c < b + c.
\]
\end{remark*}

\begin{remark*}[Predicate reading]
\[
a<b\Rightarrow a+c<b+c.
\]
\end{remark*}

\begin{remark*}[Interpretation]
Adding the same real number translates both sides by the same amount, so strict order is preserved.
\end{remark*}

\begin{dependencies}
\hyperref[def:real-ordered-field]{Ordered field},
\hyperref[thm:q-is-an-ordered-field]{ordered field axioms},
\hyperref[cor:additive-cancellation-on-q]{additive cancellation}.
\end{dependencies}

% ---------------------------------------------------------
% thm:ineq-nonstrict-add-both-sides
% ---------------------------------------------------------

\begin{theorem}[Addition Preserves Nonstrict Inequality]
\label{thm:ineq-nonstrict-add-both-sides}
For all $a, b, c \in \mathbb{R}$,
\[
  a \le b \;\Rightarrow\; a + c \le b + c.
\]
\hyperref[prf:ineq-nonstrict-add-both-sides]{\textit{Go to proof.}}
\end{theorem}



\begin{remark*}[Standard quantified statement]
\[
  \forall a, b, c \in \mathbb{R},\quad
  a \le b \;\Rightarrow\; a + c \le b + c.
\]
\end{remark*}

\begin{remark*}[Predicate reading]
\[
a\le b\Rightarrow a+c\le b+c.
\]
\end{remark*}

\begin{remark*}[Interpretation]
Nonstrict order is also stable under translation by a common real number.
\end{remark*}

\begin{dependencies}
\hyperref[def:real-ordered-field]{Ordered field},
\hyperref[thm:ineq-add-both-sides]{addition preserves strict inequality}.
\end{dependencies}

% ---------------------------------------------------------
% thm:ineq-add-inequalities
% ---------------------------------------------------------

\begin{theorem}[Addition of Strict Inequalities]
\label{thm:ineq-add-inequalities}
For all $a, b, c, d \in \mathbb{R}$,
\[
  a < b \;\land\; c < d \;\Rightarrow\; a + c < b + d.
\]
\hyperref[prf:ineq-add-inequalities]{\textit{Go to proof.}}
\end{theorem}



\begin{remark*}[Standard quantified statement]
\[
  \forall a, b, c, d \in \mathbb{R},\quad
  a < b \;\land\; c < d \;\Rightarrow\; a + c < b + d.
\]
\end{remark*}

\begin{remark*}[Predicate reading]
\[
(a<b\land c<d)\Rightarrow a+c<b+d.
\]
\end{remark*}

\begin{remark*}[Interpretation]
Strict inequalities can be added componentwise because each comparison contributes positive slack.
\end{remark*}

\begin{dependencies}
\hyperref[thm:ineq-add-both-sides]{addition preserves strict inequality},
\hyperref[thm:ineq-transitivity-strict]{transitivity of strict inequality}.
\end{dependencies}

% ---------------------------------------------------------
% thm:ineq-nonstrict-add-inequalities
% ---------------------------------------------------------

\begin{theorem}[Addition of Nonstrict Inequalities]
\label{thm:ineq-nonstrict-add-inequalities}
For all $a, b, c, d \in \mathbb{R}$,
\[
  a \le b \;\land\; c \le d \;\Rightarrow\; a + c \le b + d.
\]
\hyperref[prf:ineq-nonstrict-add-inequalities]{\textit{Go to proof.}}
\end{theorem}



\begin{remark*}[Standard quantified statement]
\[
  \forall a, b, c, d \in \mathbb{R},\quad
  a \le b \;\land\; c \le d \;\Rightarrow\; a + c \le b + d.
\]
\end{remark*}

\begin{remark*}[Predicate reading]
\[
(a\le b\land c\le d)\Rightarrow a+c\le b+d.
\]
\end{remark*}

\begin{remark*}[Interpretation]
Nonstrict inequalities add componentwise without creating more separation than the hypotheses provide.
\end{remark*}

\begin{dependencies}
\hyperref[thm:ineq-nonstrict-add-both-sides]{addition preserves nonstrict inequality}.
\end{dependencies}

\begin{lemmabox}{Lemma (Positive Sum)}
\begin{lemma}[Positive Sum]
\label{lem:positive-sum}
If $a,b\in\mathbb{R}$, $a>0$, and $b>0$, then
\[
a+b>0.
\]
\hyperref[prf:positive-sum]{Go to proof.}
\end{lemma}
\end{lemmabox}

\begin{remark*}[Standard quantified statement]
\[
\forall a,b\in\mathbb{R}\;\bigl((a>0\land b>0)\Rightarrow a+b>0\bigr).
\]
\end{remark*}

\begin{remark*}[Predicate reading]
\[
a>0\land b>0\Rightarrow a+b>0.
\]
\end{remark*}

\begin{remark*}[Interpretation]
The positive half-line is closed under addition.
\end{remark*}

\begin{dependencies}
\hyperref[thm:r-is-an-ordered-field]{$\mathbb{R}$ Is an Ordered Field};
\hyperref[thm:ineq-add-inequalities]{Addition of Strict Inequalities}.
\end{dependencies}

\begin{propositionbox}{Proposition (Subtraction of Inequalities Is Not Valid in General)}
\begin{proposition}[Subtraction of Inequalities Is Not Valid in General]
\label{prop:subtraction-of-inequalities-not-valid}
There exist real numbers $a,b,c,d$ such that $a>b$ and $c>d$, but neither
\[
a-c>b-d
\quad\text{nor}\quad
c-a>d-b
\]
holds.
\hyperref[prf:subtraction-of-inequalities-not-valid]{Go to proof.}
\end{proposition}
\end{propositionbox}

\begin{remark*}[Standard quantified statement]
\[
\exists a,b,c,d\in\mathbb{R}\;
\bigl(a>b\land c>d\land \neg(a-c>b-d)\land\neg(c-a>d-b)\bigr).
\]
\end{remark*}

\begin{remark*}[Predicate reading]
\[
a>b\land c>d\nRightarrow a-c>b-d.
\]
\end{remark*}

\begin{remark*}[Interpretation]
Inequalities may be added componentwise, but subtracting one inequality from
another has no general order rule.
\end{remark*}

\begin{dependencies}
\hyperref[thm:ineq-add-inequalities]{Addition of Strict Inequalities};
\hyperref[thm:ineq-multiply-negative]{Multiplication by a Negative Scalar Reverses Strict Inequality}.
\end{dependencies}

% ---------------------------------------------------------
% thm:ineq-mixed-add
% ---------------------------------------------------------

\begin{theorem}[Mixed Addition of Inequalities]
\label{thm:ineq-mixed-add}
For all $a, b, c, d \in \mathbb{R}$,
\[
  a \le b \;\land\; c < d \;\Rightarrow\; a + c < b + d.
\]
\hyperref[prf:ineq-mixed-add]{\textit{Go to proof.}}
\end{theorem}



\begin{remark*}[Standard quantified statement]
\[
  \forall a, b, c, d \in \mathbb{R},\quad
  a \le b \;\land\; c < d \;\Rightarrow\; a + c < b + d.
\]
\end{remark*}

\begin{remark*}[Predicate reading]
\[
(a<b\land c\le d)\Rightarrow a+c<b+d.
\]
\end{remark*}

\begin{remark*}[Interpretation]
One strict comparison and one nonstrict comparison still yield a strict comparison after addition.
\end{remark*}

\begin{dependencies}
\hyperref[thm:ineq-nonstrict-add-both-sides]{addition preserves nonstrict inequality},
\hyperref[thm:ineq-add-both-sides]{addition preserves strict inequality},
\hyperref[thm:ineq-transitivity-mixed]{mixed transitivity}.
\end{dependencies}

% ---------------------------------------------------------
% thm:ineq-multiply-positive
% ---------------------------------------------------------

\begin{theorem}[Multiplication by a Positive Scalar Preserves Strict Inequality]
\label{thm:ineq-multiply-positive}
For all $a, b, c \in \mathbb{R}$,
\[
  a < b \;\land\; 0 < c \;\Rightarrow\; ac < bc.
\]
\hyperref[prf:ineq-multiply-positive]{\textit{Go to proof.}}
\end{theorem}



\begin{remark*}[Standard quantified statement]
\[
  \forall a, b, c \in \mathbb{R},\quad
  a < b \;\land\; 0 < c \;\Rightarrow\; ac < bc.
\]
\end{remark*}

\begin{remark*}[Predicate reading]
\[
(a<b\land c>0)\Rightarrow ac<bc.
\]
\end{remark*}

\begin{remark*}[Interpretation]
Multiplication by a positive scalar rescales distances without reversing orientation.
\end{remark*}

\begin{dependencies}
\hyperref[def:real-ordered-field]{Ordered field},
multiply positive,
\hyperref[def:positive-cone]{positive cone}.
\end{dependencies}

% ---------------------------------------------------------
% thm:ineq-multiply-negative
% ---------------------------------------------------------

\begin{theorem}[Multiplication by a Negative Scalar Reverses Strict Inequality]
\label{thm:ineq-multiply-negative}
For all $a, b, c \in \mathbb{R}$,
\[
  a < b \;\land\; c < 0 \;\Rightarrow\; ac > bc.
\]
\hyperref[prf:ineq-multiply-negative]{\textit{Go to proof.}}
\end{theorem}



\begin{remark*}[Standard quantified statement]
\[
  \forall a, b, c \in \mathbb{R},\quad
  a < b \;\land\; c < 0 \;\Rightarrow\; ac > bc.
\]
\end{remark*}

\begin{remark*}[Predicate reading]
\[
(a<b\land c<0)\Rightarrow bc<ac.
\]
\end{remark*}

\begin{remark*}[Interpretation]
Multiplication by a negative scalar reverses the orientation of the real line.
\end{remark*}

\begin{dependencies}
\hyperref[def:real-ordered-field]{Ordered field},
multiply negative,
\hyperref[def:positive-cone]{positive cone}.
\end{dependencies}

% ---------------------------------------------------------
% thm:ineq-nonstrict-multiply-positive
% ---------------------------------------------------------

\begin{theorem}[Multiplication by a Positive Scalar Preserves Nonstrict Inequality]
\label{thm:ineq-nonstrict-multiply-positive}
For all $a, b, c \in \mathbb{R}$,
\[
  a \le b \;\land\; 0 < c \;\Rightarrow\; ac \le bc.
\]
\hyperref[prf:ineq-nonstrict-multiply-positive]{\textit{Go to proof.}}
\end{theorem}



\begin{remark*}[Standard quantified statement]
\[
  \forall a, b, c \in \mathbb{R},\quad
  a \le b \;\land\; 0 < c \;\Rightarrow\; ac \le bc.
\]
\end{remark*}

\begin{remark*}[Predicate reading]
\[
(a\le b\land c>0)\Rightarrow ac\le bc.
\]
\end{remark*}

\begin{remark*}[Interpretation]
Positive scalar multiplication preserves nonstrict order, including equality cases.
\end{remark*}

\begin{dependencies}
\hyperref[thm:ineq-multiply-positive]{multiplication by positive preserves strict inequality}.
\end{dependencies}

% ---------------------------------------------------------
% thm:ineq-nonstrict-multiply-nonneg
% ---------------------------------------------------------

\begin{theorem}[Multiplication by a Nonneg Scalar Preserves Nonstrict Inequality]
\label{thm:ineq-nonstrict-multiply-nonneg}
For all $a, b, c \in \mathbb{R}$,
\[
  a \le b \;\land\; 0 \le c \;\Rightarrow\; ac \le bc.
\]
\hyperref[prf:ineq-nonstrict-multiply-nonneg]{\textit{Go to proof.}}
\end{theorem}



\begin{remark*}[Standard quantified statement]
\[
  \forall a, b, c \in \mathbb{R},\quad
  a \le b \;\land\; 0 \le c \;\Rightarrow\; ac \le bc.
\]
\end{remark*}

\begin{remark*}[Predicate reading]
\[
(a\le b\land 0\le c)\Rightarrow ac\le bc.
\]
\end{remark*}

\begin{remark*}[Interpretation]
This extends
\hyperref[thm:ineq-nonstrict-multiply-positive]{the positive-scalar result}
to include $c = 0$, where both sides become $0$ and the inequality holds
trivially.
\end{remark*}

\begin{dependencies}
\hyperref[thm:ineq-nonstrict-multiply-positive]{multiplication by positive preserves nonstrict inequality},
\hyperref[def:real-ordered-field]{ordered field}.
\end{dependencies}

\begin{lemmabox}{Lemma (Positive Product)}
\begin{lemma}[Positive Product]\label{lem:positive-product}
If $a,b\in\mathbb{R}$, $a>0$, and $b>0$, then
\[
ab>0.
\]
\hyperref[prf:positive-product]{Go to proof.}
\end{lemma}
\end{lemmabox}

\begin{remark*}[Standard quantified statement]
\[
\forall a,b\in\mathbb{R},\quad (a>0\land b>0)\Rightarrow ab>0.
\]
\end{remark*}

\begin{remark*}[Predicate reading]
\[
(a>0\land b>0)\Rightarrow ab>0.
\]
\end{remark*}

\begin{remark*}[Interpretation]
The positive cone is closed under multiplication.
\end{remark*}

\begin{dependencies}
\begin{itemize}
  \item \hyperref[thm:r-is-an-ordered-field]{$\mathbb{R}$ Is an Ordered Field}
  \item \hyperref[thm:ineq-multiply-positive]{Multiplication by a Positive Scalar Preserves Strict Inequality}
\end{itemize}
\end{dependencies}

\begin{lemmabox}{Lemma (Negative Times Negative Is Positive)}
\begin{lemma}[Negative Times Negative Is Positive]\label{lem:negative-times-negative-is-positive}
If $a,b\in\mathbb{R}$, $a<0$, and $b<0$, then
\[
ab>0.
\]
\hyperref[prf:negative-times-negative-is-positive]{Go to proof.}
\end{lemma}
\end{lemmabox}

\begin{remark*}[Standard quantified statement]
\[
\forall a,b\in\mathbb{R},\quad (a<0\land b<0)\Rightarrow ab>0.
\]
\end{remark*}

\begin{remark*}[Predicate reading]
\[
(a<0\land b<0)\Rightarrow ab>0.
\]
\end{remark*}

\begin{remark*}[Interpretation]
Multiplying two negative factors reverses sign twice, producing a positive product.
\end{remark*}

\begin{dependencies}
\begin{itemize}
  \item \hyperref[thm:r-is-an-ordered-field]{$\mathbb{R}$ Is an Ordered Field}
  \item \hyperref[thm:ineq-multiply-negative]{Multiplication by a Negative Scalar Reverses Strict Inequality}
\end{itemize}
\end{dependencies}

\begin{lemmabox}{Lemma (Positive Times Negative Is Negative)}
\begin{lemma}[Positive Times Negative Is Negative]\label{lem:positive-times-negative-is-negative}
If $a,b\in\mathbb{R}$, $a>0$, and $b<0$, then
\[
ab<0.
\]
\hyperref[prf:positive-times-negative-is-negative]{Go to proof.}
\end{lemma}
\end{lemmabox}

\begin{remark*}[Standard quantified statement]
\[
\forall a,b\in\mathbb{R},\quad (a>0\land b<0)\Rightarrow ab<0.
\]
\end{remark*}

\begin{remark*}[Predicate reading]
\[
(a>0\land b<0)\Rightarrow ab<0.
\]
\end{remark*}

\begin{remark*}[Interpretation]
A positive factor preserves the negative sign of the other factor.
\end{remark*}

\begin{dependencies}
\begin{itemize}
  \item \hyperref[thm:r-is-an-ordered-field]{$\mathbb{R}$ Is an Ordered Field}
  \item \hyperref[thm:ineq-multiply-positive]{Multiplication by a Positive Scalar Preserves Strict Inequality}
\end{itemize}
\end{dependencies}

\begin{lemmabox}{Lemma (Negative Times Positive Is Negative)}
\begin{lemma}[Negative Times Positive Is Negative]\label{lem:negative-times-positive-is-negative}
If $a,b\in\mathbb{R}$, $a<0$, and $b>0$, then
\[
ab<0.
\]
\hyperref[prf:negative-times-positive-is-negative]{Go to proof.}
\end{lemma}
\end{lemmabox}

\begin{remark*}[Standard quantified statement]
\[
\forall a,b\in\mathbb{R},\quad (a<0\land b>0)\Rightarrow ab<0.
\]
\end{remark*}

\begin{remark*}[Predicate reading]
\[
(a<0\land b>0)\Rightarrow ab<0.
\]
\end{remark*}

\begin{remark*}[Interpretation]
Multiplication by a positive factor preserves strict negativity.
\end{remark*}

\begin{dependencies}
\begin{itemize}
  \item \hyperref[thm:r-is-an-ordered-field]{$\mathbb{R}$ Is an Ordered Field}
  \item \hyperref[thm:ineq-multiply-positive]{Multiplication by a Positive Scalar Preserves Strict Inequality}
\end{itemize}
\end{dependencies}

\begin{lemmabox}{Lemma (Order and Subtraction)}
\begin{lemma}[Order and Subtraction]\label{lem:order-and-subtraction}
For $a,b\in\mathbb{R}$,
\[
a<b \;\Longleftrightarrow\; b-a>0.
\]
\hyperref[prf:order-and-subtraction]{Go to proof.}
\end{lemma}
\end{lemmabox}

\begin{remark*}[Standard quantified statement]
\[
\forall a,b\in\mathbb{R},\quad a<b\Longleftrightarrow b-a>0.
\]
\end{remark*}

\begin{remark*}[Predicate reading]
\[
a<b\Longleftrightarrow b-a>0.
\]
\end{remark*}

\begin{remark*}[Interpretation]
Strict order can be tested by moving both terms to one side and checking positivity.
\end{remark*}

\begin{dependencies}
\begin{itemize}
  \item \hyperref[thm:ineq-add-both-sides]{Addition Preserves Strict Inequality}
  \item \hyperref[thm:r-is-an-ordered-field]{$\mathbb{R}$ Is an Ordered Field}
\end{itemize}
\end{dependencies}

\begin{lemmabox}{Lemma (Non-Strict Order and Subtraction)}
\begin{lemma}[Non-Strict Order and Subtraction]\label{lem:non-strict-order-and-subtraction}
For $a,b\in\mathbb{R}$,
\[
a\le b \;\Longleftrightarrow\; b-a\ge 0.
\]
\hyperref[prf:non-strict-order-and-subtraction]{Go to proof.}
\end{lemma}
\end{lemmabox}

\begin{remark*}[Standard quantified statement]
\[
\forall a,b\in\mathbb{R},\quad a\le b\Longleftrightarrow b-a\ge0.
\]
\end{remark*}

\begin{remark*}[Predicate reading]
\[
a\le b\Longleftrightarrow b-a\ge0.
\]
\end{remark*}

\begin{remark*}[Interpretation]
Nonstrict order is equivalent to nonnegativity of the difference.
\end{remark*}

\begin{dependencies}
\begin{itemize}
  \item \hyperref[thm:ineq-nonstrict-add-both-sides]{Addition Preserves Nonstrict Inequality}
  \item \hyperref[thm:r-is-an-ordered-field]{$\mathbb{R}$ Is an Ordered Field}
\end{itemize}
\end{dependencies}

\begin{lemmabox}{Lemma (Division by Positive Preserves Order)}
\begin{lemma}[Division by Positive Preserves Order]\label{lem:division-by-positive-preserves-order}
If $a,b,c\in\mathbb{R}$, $a<b$, and $c>0$, then
\[
\frac ac<\frac bc.
\]
\hyperref[prf:division-by-positive-preserves-order]{Go to proof.}
\end{lemma}
\end{lemmabox}

\begin{remark*}[Standard quantified statement]
\[
\forall a,b,c\in\mathbb{R},\quad (a<b\land c>0)\Rightarrow \frac ac<\frac bc.
\]
\end{remark*}

\begin{remark*}[Predicate reading]
\[
(a<b\land c>0)\Rightarrow a/c<b/c.
\]
\end{remark*}

\begin{remark*}[Interpretation]
Division by a positive number is multiplication by a positive reciprocal, so it preserves order.
\end{remark*}

\begin{dependencies}
\begin{itemize}
  \item \hyperref[thm:ineq-multiply-positive]{Multiplication by a Positive Scalar Preserves Strict Inequality}
  \item \hyperref[thm:r-is-an-ordered-field]{$\mathbb{R}$ Is an Ordered Field}
\end{itemize}
\end{dependencies}

\begin{lemmabox}{Lemma (Division by Negative Reverses Order)}
\begin{lemma}[Division by Negative Reverses Order]\label{lem:division-by-negative-reverses-order}
If $a,b,c\in\mathbb{R}$, $a<b$, and $c<0$, then
\[
\frac bc<\frac ac.
\]
\hyperref[prf:division-by-negative-reverses-order]{Go to proof.}
\end{lemma}
\end{lemmabox}

\begin{remark*}[Standard quantified statement]
\[
\forall a,b,c\in\mathbb{R},\quad (a<b\land c<0)\Rightarrow \frac bc<\frac ac.
\]
\end{remark*}

\begin{remark*}[Predicate reading]
\[
(a<b\land c<0)\Rightarrow b/c<a/c.
\]
\end{remark*}

\begin{remark*}[Interpretation]
Division by a negative number is multiplication by a negative reciprocal, so it reverses order.
\end{remark*}

\begin{dependencies}
\begin{itemize}
  \item \hyperref[thm:ineq-multiply-negative]{Multiplication by a Negative Scalar Reverses Strict Inequality}
  \item \hyperref[thm:r-is-an-ordered-field]{$\mathbb{R}$ Is an Ordered Field}
\end{itemize}
\end{dependencies}

\begin{lemmabox}{Lemma (Positive Multiplication Cancellation)}
\begin{lemma}[Positive Multiplication Cancellation]
\label{lem:positive-multiplication-cancellation}
If $a,b,c\in\mathbb{R}$, $c>0$, and
\[
ac<bc,
\]
then $a<b$.
\hyperref[prf:positive-multiplication-cancellation]{Go to proof.}
\end{lemma}
\end{lemmabox}

\begin{remark*}[Standard quantified statement]
\[
\forall a,b,c\in\mathbb{R}\;
\bigl((c>0\land ac<bc)\Rightarrow a<b\bigr).
\]
\end{remark*}

\begin{remark*}[Predicate reading]
\[
c>0\land ac<bc\Rightarrow a<b.
\]
\end{remark*}

\begin{remark*}[Interpretation]
A positive common factor can be cancelled from a strict inequality without
changing the order direction.
\end{remark*}

\begin{dependencies}
\hyperref[lem:division-by-positive-preserves-order]{Division by Positive Preserves Order};
\hyperref[thm:r-is-an-ordered-field]{$\mathbb{R}$ Is an Ordered Field}.
\end{dependencies}

\begin{propositionbox}{Proposition (Division of Inequalities Is Not Valid in General)}
\begin{proposition}[Division of Inequalities Is Not Valid in General]
\label{prop:division-of-inequalities-not-valid}
There exist nonzero real numbers $a,b,c,d$ such that $a>b$ and $c>d$, but
neither
\[
\frac{a}{c}>\frac{b}{d}
\quad\text{nor}\quad
\frac{c}{a}>\frac{d}{b}
\]
holds.
\hyperref[prf:division-of-inequalities-not-valid]{Go to proof.}
\end{proposition}
\end{propositionbox}

\begin{remark*}[Standard quantified statement]
\[
\exists a,b,c,d\in\mathbb{R}\;
\biggl(a,b,c,d\ne0\land a>b\land c>d
\land \neg\left(\frac ac>\frac bd\right)
\land \neg\left(\frac ca>\frac db\right)\biggr).
\]
\end{remark*}

\begin{remark*}[Predicate reading]
\[
a>b\land c>d\nRightarrow a/c>b/d.
\]
\end{remark*}

\begin{remark*}[Interpretation]
Dividing one inequality by another is not a valid order operation; division is
controlled by the sign of the divisor, not by a second comparison.
\end{remark*}

\begin{dependencies}
\hyperref[lem:division-by-positive-preserves-order]{Division by Positive Preserves Order};
\hyperref[lem:division-by-negative-reverses-order]{Division by Negative Reverses Order}.
\end{dependencies}

% ---------------------------------------------------------
% thm:ineq-squeeze
% ---------------------------------------------------------

\begin{theorembox}{Theorem (Squeeze)}
\begin{theorem}[Squeeze]
\label{thm:ineq-squeeze}
For all $a, b, c \in \mathbb{R}$,
\[
  a \le b \le c \;\land\; a = c \;\Rightarrow\; b = a.
\]
\hyperref[prf:ineq-squeeze]{\textit{Go to proof.}}
\end{theorem}
\end{theorembox}



\begin{remark*}[Standard quantified statement]
\[
  \forall a, b, c \in \mathbb{R},\quad
  a \le b \le c \;\land\; a = c \;\Rightarrow\; b = a.
\]
\end{remark*}

\begin{remark*}[Predicate reading]
\[
(a\le b\le c\land a=c)\Rightarrow b=a.
\]
\end{remark*}

\begin{remark*}[Interpretation]
The squeeze principle forces equality when a value is simultaneously
bounded above and below by equal quantities.  The continuous analogue
(if $f(x) \le g(x) \le h(x)$ and $f(x) \to L$ and $h(x) \to L$ then
$g(x) \to L$) is the central tool for computing limits.  The algebraic
version stated here is its foundation.
\end{remark*}

\begin{dependencies}
\hyperref[def:real-ordered-field]{Ordered field},
\hyperref[cor:additive-cancellation-on-q]{additive cancellation}.
\end{dependencies}

% ---------------------------------------------------------
% thm:ineq-transitivity-strict
% ---------------------------------------------------------

\begin{theorem}[Transitivity of Strict Inequality]
\label{thm:ineq-transitivity-strict}
For all $a, b, c \in \mathbb{R}$,
\[
  a < b \;\land\; b < c \;\Rightarrow\; a < c.
\]
\hyperref[prf:ineq-transitivity-strict]{\textit{Go to proof.}}
\end{theorem}



\begin{remark*}[Standard quantified statement]
\[
  \forall a, b, c \in \mathbb{R},\quad
  a < b \;\land\; b < c \;\Rightarrow\; a < c.
\]
\end{remark*}

\begin{remark*}[Predicate reading]
\[
(a<b\land b<c)\Rightarrow a<c.
\]
\end{remark*}

\begin{remark*}[Interpretation]
Strict order comparisons chain: an intermediate point can be removed without losing the outer comparison.
\end{remark*}

\begin{dependencies}
\hyperref[def:real-ordered-field]{Ordered field},
\hyperref[def:positive-cone]{positive cone}.
\end{dependencies}

% ---------------------------------------------------------
% thm:ineq-transitivity-mixed
% ---------------------------------------------------------

\begin{theorem}[Mixed Transitivity]
\label{thm:ineq-transitivity-mixed}
For all $a, b, c \in \mathbb{R}$,
\[
  a \le b \;\land\; b < c \;\Rightarrow\; a < c.
\]
\hyperref[prf:ineq-transitivity-mixed]{\textit{Go to proof.}}
\end{theorem}



\begin{remark*}[Standard quantified statement]
\[
  \forall a, b, c \in \mathbb{R},\quad
  a \le b \;\land\; b < c \;\Rightarrow\; a < c.
\]
\end{remark*}

\begin{remark*}[Predicate reading]
\[
(a<b\land b\le c)\Rightarrow a<c.
\]
\end{remark*}

\begin{remark*}[Interpretation]
A strict step followed by a nonstrict step remains strict across the endpoints.
\end{remark*}

\begin{dependencies}
\hyperref[def:real-ordered-field]{Ordered field},
\hyperref[thm:ineq-transitivity-strict]{transitivity of strict inequality}.
\end{dependencies}

% ---------------------------------------------------------
% thm:ineq-reciprocal-positive
% ---------------------------------------------------------

\begin{theorem}[Reciprocal Reverses Strict Inequality for Positives]
\label{thm:ineq-reciprocal-positive}
For all $a, b \in \mathbb{R}$,
\[
  0 < a < b \;\Rightarrow\; 0 < \frac{1}{b} < \frac{1}{a}.
\]
\hyperref[prf:ineq-reciprocal-positive]{\textit{Go to proof.}}
\end{theorem}



\begin{remark*}[Standard quantified statement]
\[
  \forall a, b \in \mathbb{R},\quad
  0 < a < b \;\Rightarrow\;
  0 < \tfrac{1}{b} < \tfrac{1}{a}.
\]
\end{remark*}

\begin{remark*}[Predicate reading]
\[
a>0\Rightarrow a^{-1}>0.
\]
\end{remark*}

\begin{remark*}[Interpretation]
The reciprocal of a positive real stays on the positive side of the ordered field.
\end{remark*}

\begin{dependencies}
\hyperref[def:real-ordered-field]{Ordered field},
\hyperref[thm:sign-of-reciprocal-r]{sign of a reciprocal on $\mathbb{R}$},
\hyperref[thm:reciprocal-reverses-order-on-positives-r]{reciprocal order for positive real elements},
\hyperref[thm:ineq-multiply-positive]{multiplication by positive preserves strict inequality}.
\end{dependencies}

% ---------------------------------------------------------
% thm:ineq-reciprocal-flip
% ---------------------------------------------------------

\begin{theorem}[Reciprocal Equivalence for Positives]
\label{thm:ineq-reciprocal-flip}
For all $a, b \in \mathbb{R}$ with $a > 0$ and $b > 0$,
\[
  a < b \iff \frac{1}{b} < \frac{1}{a}.
\]
\hyperref[prf:ineq-reciprocal-flip]{\textit{Go to proof.}}
\end{theorem}



\begin{remark*}[Standard quantified statement]
\[
  \forall a, b \in \mathbb{R},\;
  a > 0 \;\land\; b > 0 \;\Rightarrow\;
  \bigl(a < b \iff \tfrac{1}{b} < \tfrac{1}{a}\bigr).
\]
\end{remark*}

\begin{remark*}[Predicate reading]
\[
(0<a<b)\Rightarrow b^{-1}<a^{-1}.
\]
\end{remark*}

\begin{remark*}[Interpretation]
On positive inputs, reciprocation reverses order because the larger denominator gives the smaller reciprocal.
\end{remark*}

\begin{dependencies}
\hyperref[thm:ineq-reciprocal-positive]{reciprocal reverses strict inequality for positives},
\hyperref[thm:reciprocal-reverses-order-on-positives-r]{reciprocal order for positive real elements},
\hyperref[thm:reciprocal-involutive-on-r]{reciprocal is involutive on $\mathbb{R}$}.
\end{dependencies}

\begin{theorembox}{Theorem (Squares Are Nonnegative)}
\begin{theorem}[Squares Are Nonnegative]
\label{thm:square-nonnegative}
For every $a\in\mathbb{R}$,
\[
a^2\ge0.
\]
\hyperref[prf:square-nonnegative]{Go to proof.}
\end{theorem}
\end{theorembox}

\begin{remark*}[Standard quantified statement]
\[
\forall a\in\mathbb{R}\;(a^2\ge0).
\]
\end{remark*}

\begin{remark*}[Predicate reading]
\[
a^2\ge0.
\]
\end{remark*}

\begin{remark*}[Interpretation]
Squaring sends every real number into the nonnegative half-line.
\end{remark*}

\begin{dependencies}
\hyperref[lem:positive-product]{Positive Product};
\hyperref[lem:negative-times-negative-is-positive]{Negative Times Negative Is Positive};
\hyperref[thm:real-order-trichotomy]{Trichotomy for Real Order}.
\end{dependencies}

% ---------------------------------------------------------
% thm:ineq-square-monotone
% ---------------------------------------------------------

\begin{definitionbox}{Definition (Square-Root Function)}
\begin{definition}[Square-Root Function on Nonnegative Reals]
\label{def:real-square-root-function}
The \emph{square-root function} is the function
\[
\sqrt{\cdot}:[0,\infty)\to[0,\infty)
\]
defined by sending each $a\ge0$ to the unique nonnegative real number
$\sqrt a$ such that
\[
(\sqrt a)^2=a.
\]
\end{definition}
\end{definitionbox}
\begin{remark*}[Standard quantified statement]
\[
\forall a\in\mathbb{R}\;
\bigl(a\ge0\Rightarrow
(\sqrt a\ge0\land(\sqrt a)^2=a)\bigr).
\]
\end{remark*}
\begin{remark*}[Predicate reading]
\[
\sqrt{\cdot}:[0,\infty)\to[0,\infty)
\quad\text{and}\quad
(\sqrt a)^2=a.
\]
\end{remark*}
\begin{remark*}[Interpretation]
Analysis uses the square root as a function on the nonnegative real half-line;
its construction and uniqueness come from the real-number development.
\end{remark*}
\begin{dependencies}
\hyperref[def:square-root]{Square Root};
\hyperref[thm:existence-of-square-roots]{Existence of Square Roots};
\hyperref[thm:uniqueness-of-nonnegative-square-roots]{Uniqueness of Nonnegative Square Roots}.
\end{dependencies}

\begin{theorem}[Square Root Is Nonnegative]
\label{thm:sqrt-nonnegative}
For every $a\in\mathbb{R}$ with $a\ge0$,
\[
\sqrt a\ge0.
\]
\hyperref[prf:sqrt-nonnegative]{\textit{Go to proof.}}
\end{theorem}
\begin{remark*}[Standard quantified statement]
\[
\forall a\in\mathbb{R}\;(a\ge0\Rightarrow \sqrt a\ge0).
\]
\end{remark*}
\begin{remark*}[Predicate reading]
\[
a\ge0\Rightarrow \sqrt a\ge0.
\]
\end{remark*}
\begin{remark*}[Interpretation]
Square root always returns the nonnegative root.
\end{remark*}
\begin{dependencies}
\hyperref[def:real-square-root-function]{Square-Root Function on Nonnegative Reals}.
\end{dependencies}

\begin{theorem}[Square Root Squares Back]
\label{thm:sqrt-square}
For every $a\in\mathbb{R}$ with $a\ge0$,
\[
(\sqrt a)^2=a.
\]
\hyperref[prf:sqrt-square]{\textit{Go to proof.}}
\end{theorem}
\begin{remark*}[Standard quantified statement]
\[
\forall a\in\mathbb{R}\;(a\ge0\Rightarrow(\sqrt a)^2=a).
\]
\end{remark*}
\begin{remark*}[Predicate reading]
\[
a\ge0\Rightarrow(\sqrt a)^2=a.
\]
\end{remark*}
\begin{remark*}[Interpretation]
Squaring the chosen square root recovers the original nonnegative input.
\end{remark*}
\begin{dependencies}
\hyperref[def:real-square-root-function]{Square-Root Function on Nonnegative Reals}.
\end{dependencies}

\begin{theorem}[Square Roots of Zero and One]
\label{thm:sqrt-zero-one}
\[
\sqrt{0}=0
\quad\text{and}\quad
\sqrt{1}=1.
\]
\hyperref[prf:sqrt-zero-one]{\textit{Go to proof.}}
\end{theorem}
\begin{remark*}[Standard quantified statement]
\[
\sqrt{0}=0\land\sqrt{1}=1.
\]
\end{remark*}
\begin{remark*}[Predicate reading]
\[
\sqrt{0}=0,\qquad \sqrt{1}=1.
\]
\end{remark*}
\begin{remark*}[Interpretation]
The square-root function fixes the two basic nonnegative idempotent squares.
\end{remark*}
\begin{dependencies}
\hyperref[def:real-square-root-function]{Square-Root Function on Nonnegative Reals};
\hyperref[thm:uniqueness-of-nonnegative-square-roots]{Uniqueness of Nonnegative Square Roots}.
\end{dependencies}

\begin{theorem}[Square Root Is Positive on Positive Reals]
\label{thm:sqrt-positive}
For every $a\in\mathbb{R}$ with $a>0$,
\[
\sqrt a>0.
\]
\hyperref[prf:sqrt-positive]{\textit{Go to proof.}}
\end{theorem}
\begin{remark*}[Standard quantified statement]
\[
\forall a\in\mathbb{R}\;(a>0\Rightarrow\sqrt a>0).
\]
\end{remark*}
\begin{remark*}[Predicate reading]
\[
a>0\Rightarrow\sqrt a>0.
\]
\end{remark*}
\begin{remark*}[Interpretation]
A positive input cannot have square root zero, because zero squares to zero.
\end{remark*}
\begin{dependencies}
\hyperref[thm:sqrt-nonnegative]{Square Root Is Nonnegative};
\hyperref[thm:sqrt-square]{Square Root Squares Back};
\hyperref[thm:sqrt-zero-one]{Square Roots of Zero and One}.
\end{dependencies}

\begin{theorem}[Square Root of a Square]
\label{thm:sqrt-of-square-absolute-value}
For every $a\in\mathbb{R}$,
\[
\sqrt{a^2}=|a|.
\]
\hyperref[prf:sqrt-of-square-absolute-value]{\textit{Go to proof.}}
\end{theorem}
\begin{remark*}[Standard quantified statement]
\[
\forall a\in\mathbb{R}\;\sqrt{a^2}=|a|.
\]
\end{remark*}
\begin{remark*}[Predicate reading]
\[
\sqrt{a^2}=|a|.
\]
\end{remark*}
\begin{remark*}[Interpretation]
Square root reverses squaring only up to magnitude, because squaring loses
sign.
\end{remark*}
\begin{dependencies}
\hyperref[def:real-square-root-function]{Square-Root Function on Nonnegative Reals};
\hyperref[def:absolute-value]{Absolute Value};
\hyperref[thm:uniqueness-of-nonnegative-square-roots]{Uniqueness of Nonnegative Square Roots}.
\end{dependencies}

\begin{theorem}[Square Root of a Product]
\label{thm:sqrt-product}
For all $a,b\in\mathbb{R}$ with $a\ge0$ and $b\ge0$,
\[
\sqrt{ab}=\sqrt a\,\sqrt b.
\]
\hyperref[prf:sqrt-product]{\textit{Go to proof.}}
\end{theorem}
\begin{remark*}[Standard quantified statement]
\[
\forall a,b\in\mathbb{R}\;
\bigl(a\ge0\land b\ge0\Rightarrow \sqrt{ab}=\sqrt a\,\sqrt b\bigr).
\]
\end{remark*}
\begin{remark*}[Predicate reading]
\[
a,b\ge0\Rightarrow \sqrt{ab}=\sqrt a\,\sqrt b.
\]
\end{remark*}
\begin{remark*}[Interpretation]
On nonnegative real inputs, square root converts products into products.
\end{remark*}
\begin{dependencies}
\hyperref[thm:sqrt-nonnegative]{Square Root Is Nonnegative};
\hyperref[thm:sqrt-square]{Square Root Squares Back};
\hyperref[thm:uniqueness-of-nonnegative-square-roots]{Uniqueness of Nonnegative Square Roots}.
\end{dependencies}

\begin{theorem}[Square Root of a Quotient]
\label{thm:sqrt-quotient}
For all $a,b\in\mathbb{R}$ with $a\ge0$ and $b>0$,
\[
\sqrt{\frac{a}{b}}=\frac{\sqrt a}{\sqrt b}.
\]
\hyperref[prf:sqrt-quotient]{\textit{Go to proof.}}
\end{theorem}
\begin{remark*}[Standard quantified statement]
\[
\forall a,b\in\mathbb{R}\;
\left(a\ge0\land b>0\Rightarrow
\sqrt{\frac{a}{b}}=\frac{\sqrt a}{\sqrt b}\right).
\]
\end{remark*}
\begin{remark*}[Predicate reading]
\[
a\ge0\land b>0\Rightarrow \sqrt{a/b}=\sqrt a/\sqrt b.
\]
\end{remark*}
\begin{remark*}[Interpretation]
The quotient law follows from the product law and the positive denominator
keeps $\sqrt b$ nonzero.
\end{remark*}
\begin{dependencies}
\hyperref[thm:sqrt-positive]{Square Root Is Positive on Positive Reals};
\hyperref[thm:sqrt-product]{Square Root of a Product};
\hyperref[thm:sqrt-square]{Square Root Squares Back}.
\end{dependencies}

\begin{theorem}[Square Root Is Strictly Monotone on Nonnegative Reals]
\label{thm:ineq-square-root-strict-monotone}
For all $a,b\in\mathbb{R}$,
\[
0\le a<b\Rightarrow \sqrt a<\sqrt b.
\]
\hyperref[prf:ineq-square-root-strict-monotone]{\textit{Go to proof.}}
\end{theorem}
\begin{remark*}[Standard quantified statement]
\[
\forall a,b\in\mathbb{R}\;(0\le a<b\Rightarrow \sqrt a<\sqrt b).
\]
\end{remark*}
\begin{remark*}[Predicate reading]
\[
0\le a<b\Rightarrow \sqrt a<\sqrt b.
\]
\end{remark*}
\begin{remark*}[Interpretation]
Square root preserves strict order on the nonnegative half-line.
\end{remark*}
\begin{dependencies}
\hyperref[thm:ineq-square-monotone]{Square Is Strictly Monotone on Nonneg Reals};
\hyperref[thm:sqrt-square]{Square Root Squares Back}.
\end{dependencies}

\begin{lemmabox}{Lemma (Positive Powers Are Positive)}
\begin{lemma}[Positive Powers Are Positive]
\label{lem:positive-powers-are-positive}
If $x>0$ and $n\in\mathbb{N}$, then
\[
x^n>0.
\]
\hyperref[prf:positive-powers-are-positive]{Go to proof.}
\end{lemma}
\end{lemmabox}
\begin{remark*}[Standard quantified statement]
\[
\forall x\in\mathbb{R}\;\forall n\in\mathbb{N}\;
\bigl(x>0\Rightarrow x^n>0\bigr).
\]
\end{remark*}
\begin{remark*}[Predicate reading]
\[
x>0\Rightarrow x^n>0.
\]
\end{remark*}
\begin{remark*}[Negated quantified statement]
\[
\exists x\in\mathbb{R}\;\exists n\in\mathbb{N}\;
\bigl(x>0\land x^n\le 0\bigr).
\]
\end{remark*}
\begin{remark*}[Negation predicate reading]
\[
\operatorname{Positive}(x)\land\neg\operatorname{Positive}(x^n).
\]
\end{remark*}
\begin{remark*}[Interpretation]
Repeated multiplication of a positive real by itself remains positive.
\end{remark*}
\begin{dependencies}
\hyperref[def:real-ordered-field]{Ordered Field};
\hyperref[thm:ineq-multiply-positive]{Multiplication by a Positive Scalar Preserves Strict Inequality}.
\end{dependencies}

\begin{lemmabox}{Lemma (Powers Preserve Order for Positive Numbers)}
\begin{lemma}[Powers Preserve Order for Positive Numbers]
\label{lem:powers-preserve-order-for-positive-numbers}
If $0<x<y$ and $n\in\mathbb{N}$, then
\[
x^n<y^n.
\]
\hyperref[prf:powers-preserve-order-for-positive-numbers]{Go to proof.}
\end{lemma}
\end{lemmabox}
\begin{remark*}[Standard quantified statement]
\[
\forall x,y\in\mathbb{R}\;\forall n\in\mathbb{N}\;
\bigl(0<x<y\Rightarrow x^n<y^n\bigr).
\]
\end{remark*}
\begin{remark*}[Predicate reading]
\[
0<x<y\Rightarrow x^n<y^n.
\]
\end{remark*}
\begin{remark*}[Negated quantified statement]
\[
\exists x,y\in\mathbb{R}\;\exists n\in\mathbb{N}\;
\bigl(0<x<y\land y^n\le x^n\bigr).
\]
\end{remark*}
\begin{remark*}[Negation predicate reading]
\[
0<x<y\land y^n\le x^n.
\]
\end{remark*}
\begin{remark*}[Interpretation]
Integer powers are strictly increasing on the positive real axis.
\end{remark*}
\begin{dependencies}
\hyperref[lem:positive-powers-are-positive]{Positive Powers Are Positive};
\hyperref[thm:ineq-multiply-positive]{Multiplication by a Positive Scalar Preserves Strict Inequality}.
\end{dependencies}

\begin{theorem}[Square Is Strictly Monotone on Nonneg Reals]
\label{thm:ineq-square-monotone}
For all $a, b \in \mathbb{R}$,
\[
  0 \le a < b \;\Rightarrow\; a^2 < b^2.
\]
\hyperref[prf:ineq-square-monotone]{\textit{Go to proof.}}
\end{theorem}



\begin{remark*}[Standard quantified statement]
\[
  \forall a, b \in \mathbb{R},\quad
  0 \le a < b \;\Rightarrow\; a^2 < b^2.
\]
\end{remark*}

\begin{remark*}[Predicate reading]
\[
(0\le a<b)\Rightarrow a^2<b^2.
\]
\end{remark*}

\begin{remark*}[Interpretation]
Squaring is strictly increasing once inputs are restricted to the nonnegative half-line.
\end{remark*}

\begin{dependencies}
\hyperref[def:real-ordered-field]{Ordered field},
\hyperref[thm:ineq-add-inequalities]{addition of strict inequalities},
\hyperref[thm:ineq-multiply-positive]{multiplication by positive preserves strict inequality},
\hyperref[prop:product-of-positives-is-positive-q]{squares are nonnegative}.
\end{dependencies}

% ---------------------------------------------------------
% thm:ineq-square-root-monotone
% ---------------------------------------------------------

\begin{theorem}[Square Root Is Monotone on Nonneg Reals]
\label{thm:ineq-square-root-monotone}
For all $a, b \in \mathbb{R}$,
\[
  0 \le a \le b \;\Rightarrow\; \sqrt{a} \le \sqrt{b}.
\]
\hyperref[prf:ineq-square-root-monotone]{\textit{Go to proof.}}
\end{theorem}



\begin{remark*}[Standard quantified statement]
\[
  \forall a, b \in \mathbb{R},\quad
  0 \le a \le b \;\Rightarrow\; \sqrt{a} \le \sqrt{b}.
\]
\end{remark*}

\begin{remark*}[Predicate reading]
\[
(0\le a<b)\Rightarrow \sqrt{a}<\sqrt{b}.
\]
\end{remark*}

\begin{remark*}[Interpretation]
Square root is order preserving on the nonnegative half-line, including
equality cases.
\end{remark*}

\begin{dependencies}
\hyperref[def:real-square-root-function]{Square-Root Function on Nonnegative Reals};
\hyperref[thm:ineq-square-root-strict-monotone]{Square Root Is Strictly Monotone on Nonnegative Reals}.
\end{dependencies}

\begin{theorembox}{Theorem (Positive Square Comparison)}
\begin{theorem}[Positive Square Comparison]
\label{thm:positive-square-comparison}
If $a,b\in\mathbb{R}$, $a>0$, $b>0$, and
\[
a^2<b^2,
\]
then $a<b$.
\hyperref[prf:positive-square-comparison]{Go to proof.}
\end{theorem}
\end{theorembox}

\begin{remark*}[Standard quantified statement]
\[
\forall a,b\in\mathbb{R}\;
\bigl((a>0\land b>0\land a^2<b^2)\Rightarrow a<b\bigr).
\]
\end{remark*}

\begin{remark*}[Predicate reading]
\[
a,b>0\land a^2<b^2\Rightarrow a<b.
\]
\end{remark*}

\begin{remark*}[Interpretation]
Among positive real numbers, comparing squares gives the same order as
comparing the numbers themselves.
\end{remark*}

\begin{dependencies}
\hyperref[thm:ineq-square-root-monotone]{Square Root Is Monotone on Nonneg Reals};
\hyperref[thm:sqrt-of-square-absolute-value]{Square Root of a Square}.
\end{dependencies}

\begin{theorembox}{Theorem (Unit-Interval Square Bounds)}
\begin{theorem}[Unit-Interval Square Bounds]
\label{thm:unit-interval-square-bounds}
For every $a\in\mathbb{R}$,
\[
0<a<1\Rightarrow a^2<a,
\]
and
\[
a>1\Rightarrow a^2>a.
\]
\hyperref[prf:unit-interval-square-bounds]{Go to proof.}
\end{theorem}
\end{theorembox}

\begin{remark*}[Standard quantified statement]
\[
\forall a\in\mathbb{R}\;
\bigl((0<a<1\Rightarrow a^2<a)\land(a>1\Rightarrow a^2>a)\bigr).
\]
\end{remark*}

\begin{remark*}[Predicate reading]
\[
0<a<1\Rightarrow a^2<a,
\qquad
a>1\Rightarrow a^2>a.
\]
\end{remark*}

\begin{remark*}[Interpretation]
Multiplying by a positive number below one shrinks it, while multiplying a
positive number above one enlarges it.
\end{remark*}

\begin{dependencies}
\hyperref[thm:ineq-multiply-positive]{Multiplication by a Positive Scalar Preserves Strict Inequality};
\hyperref[lem:positive-multiplication-cancellation]{Positive Multiplication Cancellation}.
\end{dependencies}

% ---------------------------------------------------------
% thm:ineq-am-gm-two
% ---------------------------------------------------------

\begin{theorem}[AM-GM Inequality for Two Terms]
\label{thm:ineq-am-gm-two}
For all $a, b \ge 0$,
\[
  \sqrt{ab} \le \frac{a + b}{2}.
\]
\hyperref[prf:ineq-am-gm-two]{\textit{Go to proof.}}
\end{theorem}



\begin{remark*}[Standard quantified statement]
\[
  \forall a, b \in \mathbb{R},\;
  a \ge 0 \;\land\; b \ge 0 \;\Rightarrow\;
  \sqrt{ab} \le \frac{a + b}{2}.
\]
\end{remark*}

\begin{remark*}[Predicate reading]
\[
(a\ge0\land b\ge0)\Rightarrow \sqrt{ab}\le \frac{a+b}{2}.
\]
\end{remark*}

\begin{remark*}[Interpretation]
The arithmetic--geometric mean inequality for two terms is the simplest
case of the AM-GM family.  The standard proof proceeds by squaring both
sides (using
\hyperref[thm:ineq-square-monotone]{square monotonicity}) and reducing
to $(\sqrt{a} - \sqrt{b})^2 \ge 0$, which follows from
\hyperref[prop:product-of-positives-is-positive-q]{squares being nonnegative}.
\end{remark*}

\begin{dependencies}
\hyperref[lem:positive-product]{positive products are positive},
\hyperref[thm:ineq-square-monotone]{square is strictly monotone on nonneg reals},
\hyperref[thm:sqrt-product]{Square Root of a Product}.
\end{dependencies}

% ---------------------------------------------------------
% Named inequalities
% ---------------------------------------------------------

\begin{definitionbox}{Definition (Classical Means)}
\begin{definition}[Arithmetic, Geometric, and Harmonic Means]
\label{def:classical-means}
Let $a_1,\dots,a_n$ be positive real numbers. Their \emph{arithmetic mean},
\emph{geometric mean}, and \emph{harmonic mean} are
\[
A(a_1,\dots,a_n)=\frac{a_1+\cdots+a_n}{n},
\]
\[
G(a_1,\dots,a_n)=(a_1\cdots a_n)^{1/n},
\]
and
\[
H(a_1,\dots,a_n)=
\frac{n}{\frac1{a_1}+\cdots+\frac1{a_n}}.
\]
\end{definition}
\end{definitionbox}

\begin{remark*}[Standard quantified statement]
\[
a_1,\dots,a_n>0
\Rightarrow
A=\frac{\sum_{j=1}^{n}a_j}{n},\quad
G=\left(\prod_{j=1}^{n}a_j\right)^{1/n},\quad
H=\frac{n}{\sum_{j=1}^{n}1/a_j}.
\]
\end{remark*}

\begin{remark*}[Predicate reading]
\[
A,\;G,\;H
\]
are the additive, multiplicative, and reciprocal-average summaries of a
positive finite list.
\end{remark*}

\begin{remark*}[Interpretation]
The three classical means package the same positive data through addition,
multiplication, and reciprocals.
\end{remark*}

\begin{dependencies}
\hyperref[def:real-square-root-function]{Square-Root Function on Nonnegative Reals};
\hyperref[thm:sqrt-product]{Square Root of a Product}.
\end{dependencies}

\begin{theorembox}{Theorem (AM-GM Inequality)}
\begin{theorem}[AM-GM Inequality]
\label{thm:ineq-am-gm}
For all positive real numbers $a_1,\dots,a_n$,
\[
\frac{a_1+\cdots+a_n}{n}
\ge
(a_1\cdots a_n)^{1/n}.
\]
Equality holds if and only if $a_1=\cdots=a_n$.
\hyperref[prf:ineq-am-gm]{Go to proof.}
\end{theorem}
\end{theorembox}

\begin{remark*}[Standard quantified statement]
\[
\forall a_1,\dots,a_n>0\;
\left(
\frac{\sum_{j=1}^{n}a_j}{n}
\ge
\left(\prod_{j=1}^{n}a_j\right)^{1/n}
\right).
\]
\end{remark*}

\begin{remark*}[Predicate reading]
\[
A(a_1,\dots,a_n)\ge G(a_1,\dots,a_n).
\]
\end{remark*}

\begin{remark*}[Interpretation]
For positive data, the additive average dominates the multiplicative average.
\end{remark*}

\begin{dependencies}
\hyperref[def:classical-means]{Arithmetic, Geometric, and Harmonic Means};
\hyperref[thm:ineq-am-gm-two]{AM-GM Inequality for Two Terms};
\hyperref[thm:square-nonnegative]{Squares Are Nonnegative}.
\end{dependencies}

\begin{theorembox}{Theorem (AM-GM-HM Inequality)}
\begin{theorem}[AM-GM-HM Inequality]
\label{thm:ineq-am-gm-hm}
For all positive real numbers $a_1,\dots,a_n$,
\[
A(a_1,\dots,a_n)\ge G(a_1,\dots,a_n)\ge H(a_1,\dots,a_n).
\]
Equality holds throughout if and only if $a_1=\cdots=a_n$.
\hyperref[prf:ineq-am-gm-hm]{Go to proof.}
\end{theorem}
\end{theorembox}

\begin{remark*}[Standard quantified statement]
\[
\forall a_1,\dots,a_n>0\;
\bigl(A(a_1,\dots,a_n)\ge G(a_1,\dots,a_n)\ge H(a_1,\dots,a_n)\bigr).
\]
\end{remark*}

\begin{remark*}[Predicate reading]
\[
A\ge G\ge H.
\]
\end{remark*}

\begin{remark*}[Interpretation]
Arithmetic averaging gives the largest of the three classical means, and
harmonic averaging gives the smallest.
\end{remark*}

\begin{dependencies}
\hyperref[def:classical-means]{Arithmetic, Geometric, and Harmonic Means};
\hyperref[thm:ineq-am-gm]{AM-GM Inequality};
\hyperref[thm:ineq-reciprocal-positive]{Reciprocal order}.
\end{dependencies}

\begin{theorembox}{Theorem (Cauchy-Schwarz Inequality)}
\begin{theorem}[Cauchy-Schwarz Inequality]
\label{thm:ineq-cauchy-schwarz}
For all real numbers $a_1,\dots,a_n$ and $b_1,\dots,b_n$,
\[
\left|\sum_{j=1}^{n}a_jb_j\right|
\le
\left(\sum_{j=1}^{n}a_j^2\right)^{1/2}
\left(\sum_{j=1}^{n}b_j^2\right)^{1/2}.
\]
Equality holds if and only if one vector is a scalar multiple of the other.
\hyperref[prf:ineq-cauchy-schwarz]{Go to proof.}
\end{theorem}
\end{theorembox}

\begin{remark*}[Standard quantified statement]
\[
\forall a_1,\dots,a_n,b_1,\dots,b_n\in\mathbb{R}\;
\left|\sum_{j=1}^{n}a_jb_j\right|
\le
\left(\sum_{j=1}^{n}a_j^2\right)^{1/2}
\left(\sum_{j=1}^{n}b_j^2\right)^{1/2}.
\]
\end{remark*}

\begin{remark*}[Predicate reading]
\[
|\langle a,b\rangle|\le \|a\|_2\|b\|_2.
\]
\end{remark*}

\begin{remark*}[Interpretation]
The dot product of two finite real lists is controlled by the product of
their Euclidean lengths.
\end{remark*}

\begin{dependencies}
\hyperref[thm:square-nonnegative]{Squares Are Nonnegative};
\hyperref[thm:sqrt-product]{Square Root of a Product};
\hyperref[thm:triangle-inequality]{Triangle Inequality}.
\end{dependencies}

\begin{theorembox}{Theorem (Chebyshev's Inequality)}
\begin{theorem}[Chebyshev's Sum Inequality]
\label{thm:ineq-chebyshev-sum}
Let $a_1\le\cdots\le a_n$ and $b_1\le\cdots\le b_n$ be real sequences.
Then
\[
\left(\frac{a_1+\cdots+a_n}{n}\right)
\left(\frac{b_1+\cdots+b_n}{n}\right)
\le
\frac{a_1b_1+\cdots+a_nb_n}{n}.
\]
The same conclusion holds when both sequences are nonincreasing.
\hyperref[prf:ineq-chebyshev-sum]{Go to proof.}
\end{theorem}
\end{theorembox}

\begin{remark*}[Standard quantified statement]
\[
\bigl(a_1\le\cdots\le a_n\land b_1\le\cdots\le b_n\bigr)
\Rightarrow
\left(\frac{\sum a_j}{n}\right)\left(\frac{\sum b_j}{n}\right)
\le
\frac{\sum a_jb_j}{n}.
\]
\end{remark*}

\begin{remark*}[Predicate reading]
\[
\left(\frac{\sum a_j}{n}\right)\left(\frac{\sum b_j}{n}\right)
\le
\frac{\sum a_jb_j}{n}.
\]
\end{remark*}

\begin{remark*}[Interpretation]
When two lists move in the same order, pairing corresponding terms raises the
average product.
\end{remark*}

\begin{dependencies}
\hyperref[thm:ineq-add-inequalities]{Addition of Strict Inequalities};
\hyperref[thm:ineq-nonstrict-multiply-nonneg]{Multiplication by a Nonneg Scalar Preserves Nonstrict Inequality}.
\end{dependencies}

\begin{theorembox}{Theorem (Rearrangement Inequality)}
\begin{theorem}[Rearrangement Inequality]
\label{thm:ineq-rearrangement}
Let
\[
a_1\le\cdots\le a_n
\quad\text{and}\quad
b_1\le\cdots\le b_n.
\]
For every permutation $(b'_1,\dots,b'_n)$ of $(b_1,\dots,b_n)$,
\[
\sum_{j=1}^{n}a_jb_{n+1-j}
\le
\sum_{j=1}^{n}a_jb'_j
\le
\sum_{j=1}^{n}a_jb_j.
\]
\hyperref[prf:ineq-rearrangement]{Go to proof.}
\end{theorem}
\end{theorembox}

\begin{remark*}[Standard quantified statement]
\[
\forall \sigma\in S_n\;
\sum_{j=1}^{n}a_jb_{n+1-j}
\le
\sum_{j=1}^{n}a_jb_{\sigma(j)}
\le
\sum_{j=1}^{n}a_jb_j.
\]
\end{remark*}

\begin{remark*}[Predicate reading]
\[
\text{same-order pairing maximizes the sum of products; opposite-order pairing minimizes it.}
\]
\end{remark*}

\begin{remark*}[Interpretation]
Sorted lists produce the largest dot product when paired in the same order
and the smallest when paired in opposite order.
\end{remark*}

\begin{dependencies}
\hyperref[thm:ineq-chebyshev-sum]{Chebyshev's Sum Inequality};
\hyperref[thm:ineq-add-inequalities]{Addition of Strict Inequalities}.
\end{dependencies}

\begin{theorembox}{Theorem (Holder's Inequality)}
\begin{theorem}[Holder's Inequality]
\label{thm:ineq-holder}
Let $p,q>1$ satisfy
\[
\frac1p+\frac1q=1.
\]
For real numbers $a_1,\dots,a_n$ and $b_1,\dots,b_n$,
\[
\sum_{j=1}^{n}|a_jb_j|
\le
\left(\sum_{j=1}^{n}|a_j|^p\right)^{1/p}
\left(\sum_{j=1}^{n}|b_j|^q\right)^{1/q}.
\]
\hyperref[prf:ineq-holder]{Go to proof.}
\end{theorem}
\end{theorembox}

\begin{remark*}[Standard quantified statement]
\[
p,q>1\land p^{-1}+q^{-1}=1
\Rightarrow
\sum_{j=1}^{n}|a_jb_j|
\le
\left(\sum_{j=1}^{n}|a_j|^p\right)^{1/p}
\left(\sum_{j=1}^{n}|b_j|^q\right)^{1/q}.
\]
\end{remark*}

\begin{remark*}[Predicate reading]
\[
\|ab\|_1\le \|a\|_p\|b\|_q.
\]
\end{remark*}

\begin{remark*}[Interpretation]
Holder's inequality is the conjugate-exponent generalization of
Cauchy-Schwarz.
\end{remark*}

\begin{dependencies}
\hyperref[thm:ineq-am-gm]{AM-GM Inequality};
\hyperref[thm:ineq-cauchy-schwarz]{Cauchy-Schwarz Inequality};
\hyperref[def:absolute-value]{Absolute Value}.
\end{dependencies}

\begin{theorembox}{Theorem (Minkowski's Inequality)}
\begin{theorem}[Minkowski's Inequality]
\label{thm:ineq-minkowski}
Let $p\ge1$. For real numbers $a_1,\dots,a_n$ and $b_1,\dots,b_n$,
\[
\left(\sum_{j=1}^{n}|a_j+b_j|^p\right)^{1/p}
\le
\left(\sum_{j=1}^{n}|a_j|^p\right)^{1/p}
+
\left(\sum_{j=1}^{n}|b_j|^p\right)^{1/p}.
\]
\hyperref[prf:ineq-minkowski]{Go to proof.}
\end{theorem}
\end{theorembox}

\begin{remark*}[Standard quantified statement]
\[
p\ge1
\Rightarrow
\|a+b\|_p\le \|a\|_p+\|b\|_p.
\]
\end{remark*}

\begin{remark*}[Predicate reading]
\[
\text{finite } \ell^p\text{ length satisfies the triangle inequality.}
\]
\end{remark*}

\begin{remark*}[Interpretation]
Minkowski's inequality is the triangle inequality for finite $\ell^p$ norms.
\end{remark*}

\begin{dependencies}
\hyperref[thm:ineq-holder]{Holder's Inequality};
\hyperref[thm:triangle-inequality]{Triangle Inequality};
\hyperref[def:absolute-value]{Absolute Value}.
\end{dependencies}

% ---------------------------------------------------------
% thm:ineq-bernoulli
% ---------------------------------------------------------

\begin{theorembox}{Theorem (Bernoulli's Inequality)}
\begin{theorem}[Bernoulli's Inequality]
\label{thm:ineq-bernoulli}
For all $x \in \mathbb{R}$ with $x \ge -1$ and all $n \in \mathbb{N}$,
\[
  (1 + x)^n \ge 1 + nx.
\]
\hyperref[prf:ineq-bernoulli]{\textit{Go to proof.}}
\end{theorem}
\end{theorembox}



\begin{remark*}[Standard quantified statement]
\[
  \forall x \in \mathbb{R},\;
  \forall n \in \mathbb{N},\;
  x \ge -1 \;\Rightarrow\;
  (1 + x)^n \ge 1 + nx.
\]
\end{remark*}

\begin{remark*}[Predicate reading]
\[
(x\ge -1\land n\in\mathbb{N})\Rightarrow (1+x)^n\ge 1+nx.
\]
\end{remark*}

\begin{remark*}[Interpretation]
Bernoulli's inequality is proved by induction and is a key estimate in
analysis: it provides a linear lower bound for exponential growth and
is used to establish the divergence of the harmonic series, bounds on
compound interest, and properties of the exponential function.
\end{remark*}

\begin{dependencies}
\hyperref[def:real-ordered-field]{Ordered field},
\hyperref[thm:ineq-nonstrict-multiply-nonneg]{multiplication by nonneg preserves nonstrict inequality},
\hyperref[thm:ineq-nonstrict-add-both-sides]{addition preserves nonstrict inequality}.
\end{dependencies}
